\chapter{Introduction}
\label{chap:introduction}

\section{Problem Statement}
\label{sec:problem-statement}

\begin{itemize}
    \item GPUs have become critical resources because of the rapid growth of artificial intelligence workloads, the expansion of data centres, and the associated infrastructure and energy costs.
    \item AI workloads have heterogeneous compute, memory, latency, and throughput requirements; consequently, they cannot all be served efficiently using the same GPU allocation strategy.
    \item GPUs are often underutilized, which leads to wasted computing capacity, energy, and infrastructure investment.
    \item Kubernetes has become a major runtime and orchestration platform for containerized AI workloads.
    \item Conventional Kubernetes device allocation handles accelerators in a relatively rigid and static manner, limiting the ability to adapt GPU resources to workload-specific requirements.
\end{itemize}

\section{Objectives}
\label{sec:objectives}

\begin{itemize}
    \item Increase effective GPU utilization while preserving acceptable inference performance and workload isolation.
    \item Design and implement proofs of concept to evaluate NVIDIA Multi-Instance GPU (MIG) and Kubernetes Dynamic Resource Allocation (DRA) under representative AI inference workloads.
\end{itemize}

\section{Thesis Structure}
\label{sec:thesis-structure}

\begin{itemize}
    \item Chapter~\ref{chap:background} introduces GPU-sharing techniques, NVIDIA MIG, GPU support in Kubernetes, DRA, the Linux GPU software stack, and the NVIDIA GPU Operator.
    \item Chapter~\ref{chap:implementation} describes the experimental testbed and the two evaluated use cases: concurrent Ollama inference services and Kueue-managed batch workloads.
    \item Chapter~\ref{chap:conclusions} summarizes the main findings and identifies directions for future work.
\end{itemize}
